\documentclass[11pt,a4paper]{article}
\usepackage[margin=1in]{geometry}
\title{DATA\_EXTRACTION\_STRUCTURE\_2026-03-08\\Scientific Technical Report}
\author{Science-Chief}
\date{2026-03-08}
\begin{document}
\maketitle

\section*{1. Scope and Research Purpose}
This document explains, in paper style, what data were extracted, from which on-chain sources, why they were extracted, and how the extraction pipeline is structured for ERC8004 analysis.
Research objective: characterize identity adoption, feedback activation, concentration, and semantic usage in the ERC8004 ecosystem.

\section*{2. Data Provenance (On-chain Registries)}
Extraction was performed by interacting with Ethereum mainnet event logs from two registries:
\begin{itemize}
\item Identity Registry: \texttt{0x8004A169FB4a3325136EB29fA0ceB6D2e539a432}
\item Reputation Registry: \texttt{0x8004BAa17C55a88189AE136b182e5fdA19dE9b63}
\end{itemize}
Reference window (validated run): block range approximately \texttt{24339925} to \texttt{24605141}.

\section*{3. Why These Data Are Extracted}
The extracted events allow us to study:
\begin{itemize}
\item Identity creation dynamics (how many agents are registered and when);
\item Trust/feedback activation (how many identities receive feedback and after how long);
\item Network concentration (who gives feedback and who receives it);
\item Semantic layer (how \texttt{tag1/tag2} are used in feedback events).
\end{itemize}

\section*{4. Extraction Procedure (What We Actually Did)}
\begin{enumerate}
\item Preflight checks (configuration and run safety).
\item Query logs from both registry addresses above.
\item Decode event payloads and persist them into CSV artifacts.
\item Store run artifacts under run-scoped folders and update latest pointer.
\item Validate row counts, unique entities, and block coverage before analysis.
\end{enumerate}

\section*{5. Dataset Structure and Schema}
Main extracted files and columns:
\begin{itemize}
\item \texttt{identity\_registered.csv} (event \texttt{Registered}):
\texttt{blockNumber, transactionHash, logIndex, address, event, agentId, owner, agentURI}
\item \texttt{identity\_transfer.csv} (event \texttt{Transfer}):
\texttt{blockNumber, transactionHash, logIndex, address, event, from, to, tokenId}
\item \texttt{reputation\_newfeedback.csv} (event \texttt{NewFeedback}):
\texttt{blockNumber, transactionHash, logIndex, address, event, agentId, clientAddress, indexedTag1, feedbackIndex, value, valueDecimals, tag1, tag2, endpoint, feedbackURI, feedbackHash}
\item \texttt{identity\_metadataset.csv}, \texttt{reputation\_feedbackrevoked.csv}: currently empty in this snapshot.
\end{itemize}

Each row corresponds to one on-chain event log; columns \texttt{address/event} tie the row to the originating registry and event type.

\section*{6. Audited Quantities (Current Snapshot)}
Registered agents: 28,380; feedback events: 2,776; agents with feedback: 1,470; transfer events: 37,293.
These counts are consistent with current milestone analyses.

\section*{7. Limitations and Interpretation Guardrails}
\begin{itemize}
\item ETH-only pass: no cross-chain generalization yet.
\item Some semantic outputs depend on normalization (e.g., \texttt{tag1} heterogeneity).
\item Time interpretations may require timestamp enrichment when block-time proxy is insufficient.
\end{itemize}

\section*{8. Decision and Next Steps}
\textbf{Decision: GO for internal scientific use; HOLD for ecosystem-wide claims.}
Next: timestamp enrichment defaults, multi-chain extension, and formal semantic QA.

\section*{Appendix: Source Paths}
\textbf{Tables}\\
\begin{itemize}
\item \texttt{workspaces/science-chief/outbox/Reports/DATA\_EXTRACTION\_STRUCTURE\_2026-03-08/tables/data\_audit\_table.tex}
\item \texttt{workspaces/science-chief/outbox/Reports/DATA\_EXTRACTION\_STRUCTURE\_2026-03-08/tables/paths\_table.tex}
\end{itemize}

\end{document}
